\documentclass[11pt,a4paper]{article}
\usepackage{amsmath,amssymb,graphicx,geometry,booktabs,hyperref}
\usepackage{setspace}
\onehalfspacing
\emergencystretch=1.5em
\geometry{margin=1in}

\title{\textbf{The Benign and the Lethal Barrier:\\
A Same-Calibration Comparison of a Double Knock-Out Hedge\\
and the KIKO}}
\author{\textbf{Jihong Park} \\ Pusan National University}
\date{July 2026}

\begin{document}

\maketitle

\begin{abstract}
Two barrier structures frame Korea's corporate-hedging experience: the knock-in knock-out (KIKO) contracts whose 2008 detonation pushed operationally profitable exporters into distress, and the double knock-out option -- the premium-reducing structure studied throughout a four-paper research program on a crude importer's WTI--currency exposure. This note prices both families under one calibration and asks what actually separates them, since the surface diagnosis (``barriers are dangerous'') indicts both. The barrier mortality is nearly symmetric: the program's own budget paper measures an $89.25\%$ stress-conditional knock-out probability for its double-KO structure, and the KIKO protection leg dies in $90$--$100\%$ of the scenarios in which it is needed. What differs is everything else. The double-KO holder \emph{owns} the optionality, so the barrier extinguishes a prepaid asset and the worst case is the premium; the KIKO holder has \emph{written} levered optionality, so the barrier awakens a liability and the worst case is unbounded -- at a fifty-percent depreciation, an exporter overhedged at the documented $2.5\times$ loses roughly forty times the price of the vanilla insurance it declined to buy. And the institutional treatment differs: the program's framework measured its structure's stress mortality against a $6.16\%$ break-even and \emph{rejected its own product}, prescribing the vanilla book under stress, while KIKO's buyers -- facing an instrument whose fair financing ratio was $0.075$ calls per put, sold at $2.0$ -- had no framework at all. The conclusion is a suitability principle with teeth: the dangerous object is not the barrier but the unmeasured barrier, and the decisive asymmetry is not payoff engineering but the presence of a framework willing to fail its own trade.
\end{abstract}

\noindent\textbf{Keywords:} barrier options; knock-in; knock-out; KIKO; zero-cost structures; corporate hedging; suitability.\\
\textbf{JEL codes:} G32, G13, G28.

\section{Introduction}
\label{sec:intro}

Barrier clauses entered Korean corporate treasuries from two directions. From one came the knock-in knock-out: a ``zero-cost'' export hedge -- one purchased put extinguished by a lower barrier, two written calls activated by an upper one -- whose 2008 mass knock-in produced losses large enough to fail hundreds of profitable firms. From the other comes the structure studied across the companion research program (Park, 2026a--c): a double knock-out quanto option bought by a crude importer to cap its joint oil-currency cost, its premium reduced by barriers that extinguish the option if oil trades outside a corridor. A regulator, an auditor, or a treasurer reading the KIKO literature could reasonably conclude that barrier structures as a class are unsuitable for corporates. This note tests that conclusion by putting both structures under the same lens -- the program's calibration, pricing machinery, and survival-haircut logic -- and finds it wrong in an instructive way: on the dimension the episode made famous, \emph{hedge mortality at the barrier}, the two structures are nearly twins. The separation lies on three other axes -- who owns the optionality, what the barrier does to the owner's balance sheet, and whether any framework in the room was capable of rejecting the trade.

Section~\ref{sec:anatomy} sets out the two anatomies. Section~\ref{sec:ledger} prices an exporter's menu -- forward, vanilla put, knocked-out put, KIKO -- under one calibration. Section~\ref{sec:twin} brings in the program's double-KO quanto and its own framework's verdict on it. Section~\ref{sec:asym} isolates the two asymmetries that matter. Section~\ref{sec:conclusion} states the resulting suitability principle. All KIKO-side quantities are computed in the companion note (Park, 2026e) with the program's first-passage Monte Carlo (200{,}000 paths, 260 steps/year, Brownian-bridge barrier correction) at spot $S_0=1{,}540.64$, volatility $9.258\%$, $r_{USD}=4.0\%$, $r_{KRW}=3.5\%$, one-year horizon; program-side quantities are as published in the companion papers.

\section{Two Anatomies}
\label{sec:anatomy}

The instruments are near-mirror images (Table~\ref{tab:anatomy}). The double-KO buyer pays a premium today for an option it owns; the barriers cheapen the premium by capping the \emph{seller's} liability, and the worst outcome the buyer can experience -- knock-out included -- is the loss of what it already paid. The KIKO buyer pays nothing today and owns almost nothing (a put priced at $3.01$ KRW per USD against a vanilla value of $71.45$); it has instead \emph{written} two calls whose knock-in barrier caps nothing and creates, upon touch, a levered obligation. In option language: one structure is long convexity with a bounded left tail, the other is short convexity with an unbounded one. In balance-sheet language: one barrier extinguishes an asset, the other ignites a liability.

\begin{table}[h]
\centering
\small
\begin{tabular}{lll}
\toprule
 & Double-KO quanto (program) & KIKO (episode) \\
\midrule
Firm's option position & long 1$\times$ call & long 1$\times$ put, \textbf{short 2$\times$ calls} \\
Premium at inception & paid, visible (17{,}132 KRW/unit) & ``zero'' -- hidden transfer $\approx5\%$ of notional \\
Barrier acts on & the firm's \emph{asset} & the firm's \emph{liability} \\
Barrier event means & prepaid option dies & levered short awakens \\
Worst case & $-$premium (bounded) & $-2N(S_T-K)$ (unbounded) \\
Convexity held by firm & positive & negative \\
Fair vs sold financing & n/a (priced trade) & fair $0.075$ calls/put, sold $2.0$ \\
\bottomrule
\end{tabular}
\caption{Anatomies. Program figures from Park (2026c); KIKO figures from the companion reconstruction (Park, 2026e) at the sold geometry ($K=F+2\%$, $U=+7\%$, $L=-5\%$).}
\label{tab:anatomy}
\end{table}

\section{One Calibration, One Menu}
\label{sec:ledger}

Table~\ref{tab:menu} prices the exporter's realistic menu under the single calibration above. Two readings matter. First, the knocked-out put -- the ``honest'' barrier trade, in which the firm pays the (tiny) fair price for barrier-cheapened protection -- is a defensible instrument in a way that has nothing to do with its mortality: its worst case is $3.01$ KRW per USD, the premium. A treasurer who buys it and watches it knock out has lost lunch money, knowingly. Second, the KIKO's catastrophic column is not produced by the barrier at all but by the \emph{leverage behind} the barrier: at a $+25\%$ depreciation the $2.5\times$-overhedged firm loses $1{,}426$ KRW per dollar of flow, roughly \textbf{twenty times} the full vanilla insurance premium it was never asked to pay -- and at $+50\%$, roughly forty times. The barrier merely schedules the detonation; the warhead is the written, levered call.

\begin{table}[h]
\centering
\small
\begin{tabular}{lrrrl}
\toprule
Structure & Cost today & Worst case & Mortality when needed & Firm's net $\Delta$ at inception \\
 & (KRW/USD) & (KRW/USD flow) & & (flow $+1$ included) \\
\midrule
Unhedged & 0 & unbounded (down) & -- & $+1$ \\
Forward at $F$ & 0 & opportunity only & none & $0$ \\
Vanilla put ($K{=}F{+}2\%$) & 71.45 & $-71.45$ & none & $+1\!\rightarrow\!0$ (kinked) \\
Knocked-out put & 3.01 & $-3.01$ & $0.90$--$1.00$ & $\approx+1$ (protection thin) \\
KIKO ($n{=}2$, ``zero-cost'') & 0 (hidden $76.8$) & unbounded (up) & $0.95$--$1.00$ & $+0.18$, $\to-1$ past $U$ \\
\bottomrule
\end{tabular}
\caption{The exporter's menu under one calibration. ``Mortality when needed'': probability the protection leg is already extinguished conditional on the terminal spot ending in its coverage zone ($S_T\le0.96$--$0.98\,S_0$). KIKO's hidden cost is the inception value transfer $2C^{KI}-P^{KO}$.}
\label{tab:menu}
\end{table}

\section{The Long Twin, and a Framework That Failed Its Own Trade}
\label{sec:twin}

Set beside this menu, the program's double-KO quanto is the long twin of the knocked-out put, scaled up: a premium of KRW 17{,}132.47 per unit (KRW 34.3 billion on the full two-million-barrel position), a European knock-out rate of $43.7\%$ over the option's life, and -- the number that matters -- a \emph{stress-conditional} knock-out probability of $89.25\%$, measured by re-simulation from the stress state in the budget paper. Its hedge dies under stress about as reliably as the KIKO's put leg does. If mortality were the criterion, the two structures would hang together.

What actually happened next is the point of this note. The budget paper carries a survival haircut: a barrier discount is worth taking only if mortality-when-needed stays below a break-even it derives at $6.16\%$. Measured mortality exceeded the break-even by a factor of fourteen, and the framework did what frameworks are for -- \emph{it rejected its own structure}, prescribing that under stress the WTI book be held in the vanilla instrument, and constructing a mixed vanilla/knock-out program as the implementable compromise (Park, 2026a, \S7--8). The trade that a desk had built an entire pricing and accounting apparatus around failed its own suitability test, in print. Nothing analogous existed on the other side of the KIKO trade: a structure whose fair financing ratio was $0.075$ written calls per put was sold at $2.0$, a five-percent-of-notional inception transfer wearing a zero, to counterparties who could not price a barrier, could not observe their own delta of $-0.82$ per notional dollar, and could not have named the $0.968$ probability that a five-percent depreciation would spring the knock-in.

\section{The Two Asymmetries That Matter}
\label{sec:asym}

\paragraph{Sign and leverage, not barriers.} Both structures monetize the same actuarial fact -- paths that touch a barrier are cheap to reinsure -- and both hedges die at their barriers with high probability under stress. The economic difference is the direction of the residual claim: long-barrier structures degrade into \emph{self-insurance} (the firm reverts to unhedged, minus a known premium), while short-levered-barrier structures degrade into \emph{reverse insurance} (the firm becomes the underwriter of the disaster it feared, at $2\times$ size). The severity ordering follows from the payoff algebra, not from any subtlety of barrier pricing: bounded-below versus unbounded-below.

\paragraph{Measurement, not sophistication.} The deeper asymmetry is institutional. The program's structure was, by the standards of the KIKO literature, the more exotic object -- two assets, a quanto coupling, jump diffusion, path-dependent barriers -- yet it circulated inside an apparatus that could price it three ways, certify the engines against each other, measure its Greeks' failure modes, book it under two designation architectures, and, decisively, compute the one number ($89.25\%$ against $6.16\%$) that disqualified it. The KIKO was the simpler payoff and the more lethal holding, because it circulated with no apparatus at all. Suitability, on this evidence, is not a property of instruments but of instrument-plus-institution pairs: the same barrier clause is benign in a hand that can measure it and lethal in a hand that cannot.

\section{Conclusion}
\label{sec:conclusion}

The comparison yields one principle and one corollary. The principle: \emph{a barrier discount is a loan against the states in which you need the hedge; taking it is defensible exactly when the borrower can state the repayment terms} -- the mortality-when-needed, the post-event exposure, the worst case -- \emph{and remains solvent under them}. The double-KO quanto satisfies the solvency clause by construction (bounded loss) and, within the program, satisfied the statement clause by measurement -- and was still, correctly, benched under stress by its own framework. The KIKO failed every clause at once: unbounded downside, unmeasured mortality, and a hidden price. The corollary is the regulatory line the companion note derives from the disclosure paper: the screen that separates the two cases \emph{ex ante} is not a product ban but a quantified ratio -- derivative notional to underlying flow, marked against its barriers -- because that single disclosure prices the leverage that turns a barrier clause from a discount into a detonator.

\section{Limitations}
\label{sec:limits}

The comparison inherits the companion note's stylizations: single-settlement KIKO against the episode's multi-window strips, the program's calibration in place of 2007 vols, no credit or margin-call channel, and an exporter flow held fixed. The two protagonist structures are additionally not symmetric objects -- one is a researched design priced at fair value, the other a sold product reconstructed at its marketed geometry -- so the institutional contrast of Section~\ref{sec:twin} compares a framework's treatment of its instrument with a market's treatment of its customers, which is precisely the point but should not be mistaken for a like-for-like product test.

\begin{thebibliography}{9}

\bibitem{broadie1996} Broadie, M., and P. Glasserman (1996). ``Estimating Security Price Derivatives Using Simulation.'' \emph{Management Science} 42(2), 269--285.

\bibitem{garman1983} Garman, M. B., and S. W. Kohlhagen (1983). ``Foreign Currency Option Values.'' \emph{Journal of International Money and Finance} 2(3), 231--237.

\bibitem{park2026a} Park, J. (2026a). ``Optimal WTI--FX Hedge Ratios Under a Fixed Budget.'' Working paper, Pusan National University.

\bibitem{park2026b} Park, J. (2026b). ``IFRS 9 Cash-Flow-Hedge Accounting for a Quanto Knock-Out: Combined Versus Split Designation.'' Working paper, Pusan National University.

\bibitem{park2026c} Park, J. (2026c). ``Covariance-Aware Delta Hedging of a Quanto Knock-Out.'' Working paper, Pusan National University.

\bibitem{park2026d} Park, J. (2026d). ``Does Mandatory ESG Disclosure Change Corporate Hedging?'' Working paper, Pusan National University.

\bibitem{park2026e} Park, J. (2026e). ``KIKO Through the Program: What a Hedging Research Framework Says About Korea's 2008 Knock-In Knock-Out Disaster.'' Working paper, Pusan National University.

\bibitem{stulz1996} Stulz, R. M. (1996). ``Rethinking Risk Management.'' \emph{Journal of Applied Corporate Finance} 9(3), 8--25.

\end{thebibliography}

\end{document}
